\chapter{Orbital Transfer Basics}

This chapter reviews the energy budget of a spacecraft that moves between two circular orbits around a single attracting body.

\section{2.1 Impulsive Manoeuvres}

A velocity change $\Delta v$ applied at the periapsis raises the apoapsis radius $r_A$ while the periapsis radius $r_P$ stays where it was.

where  $\mu$  is the gravitational parameter of the central body and  $v_c = \sqrt{\mu/r}$  is the local circular speed. For a small burn, it is

\begin{equation}
\Delta v \simeq \frac{\mu}{2 r^{2} v_c}\,\Delta a
\end{equation}

and

The estimate agrees with the energy argument of Section~\ref{sec:energy} and with the tabulated cases collected in \cite{tablo2013}.

\textbf{Remark:} The same reasoning holds for a burn at the apoapsis, provided that $\Delta v$ remains small compared with $v_c$.

\begin{itemize}
  \item the thrust points along the local velocity $\vec{v}$ during the whole burn;
  \item the propellant mass $m_p$ is negligible with respect to the dry mass $m_0$;
  \item the burn time $\tau$ is much shorter than the orbital period $T$.
\end{itemize}

\begin{figure}[H]
  \centering
  \includegraphics[width=0.8\textwidth]{transfer_sketch.png}
\end{figure}

Figure 2.1: Sketch of a two-impulse transfer between coplanar orbits.

Motus omnis corporis circa Solem in orbe fit, et tempora periodica sunt in ratione sesquialtera distantiarum mediarum, quod experientia confirmat.

The total cost of the manoeuvre is therefore  $\Delta v_{\rm tot} = \Delta v_1 + \Delta v_2$ , which is smallest when the transfer ellipse touches both circles.

where

$\Delta v_1$  and  $\Delta v_2$  denote the first and the second impulse, both measured in km/s.

For a radius ratio $r_2/r_1$ above a critical value the three-impulse route becomes cheaper, and the transfer time $t_{12}$ grows accordingly.
